% Appendix B -- Kode-dokumentation

\chapter{Kode-dokumentation}
\label{AppendixB}

Dette bilag samler den dokumentation, der knytter sig direkte til den
software og de konfigurationer, som er udviklet eller ændret i
forbindelse med projektet. Hovedteksten henviser til de relevante
afsnit i dette bilag i stedet for at gengive kode og konfigurationer
i fuld længde.

% ---------------------------------------------------------------------
\section{Manager-noden}
\label{sec:appB-manager}

% TODO: Beskrivelse af wattson_servo_manager_node:
%  - hvilke ROS-topics den lytter på og publicerer
%  - hvilke services den eksponerer
%  - hvordan JSON-konfigurationen indlæses ved opstart
%  - hvordan kommandoer videresendes til ESP32-kontrolleren
% Indsæt evt. relevant kodeuddrag i \begin{verbatim} ... \end{verbatim}.

\textit{Indhold tilføjes -- se README og kildekode i workspace.}

% ---------------------------------------------------------------------
\section{JSON-konfiguration: servo\_arm\_left\_params.json}
\label{sec:appB-json}

% TODO: Beskrivelse af JSON-strukturen:
%  - hvilke nøgler hver servo-post har (servo_id, gear_ratio, gain_P,
%    gain_I, gain_D, min/max, default, feedback_enabled)
%  - hvilke værdier der er ændret i projektet og hvorfor
%  - eksempel-snippet med én servo-post

\textit{Indhold tilføjes.}

% ---------------------------------------------------------------------
\section{Driver: DriverWaveshare}
\label{sec:appB-driver}

% TODO: Kort beskrivelse af driver-klassen:
%  - hvilken rolle den spiller mellem manager-noden og ESP32
%  - hvilke metoder der er væsentlige (særligt instans-defaults
%    for feedback_enabled, jf. kapitel~\ref{ch:idrift})

\textit{Indhold tilføjes.}

% ---------------------------------------------------------------------
\section{ServoControl.compute\_command og feedback\_enabled-gating}
\label{sec:appB-feedback}

% TODO: Dokumentation af den tre-niveau-gating, der er beskrevet i
% kapitel~\ref{ch:idrift}:
%  1) feedback_enabled i JSON-konfigurationen
%  2) instans-default i DriverWaveshare
%  3) effekt i ServoControl.compute_command (cachet position vs. målt)

\textit{Indhold tilføjes.}

% ---------------------------------------------------------------------
\section{step\_response\_node}
\label{sec:appB-step}

% TODO: Beskrivelse af step_response_node:
%  - formål (logging af feedback-vinkel ved fast step)
%  - hvilke topics den lytter på/publicerer
%  - output-format (CSV) og parametre (step-amplitude, varighed)

\textit{Indhold tilføjes.}

% ---------------------------------------------------------------------
\section{Launcher-GUI}
\label{sec:appB-gui}

% TODO: Dokumentation af launcher-GUI'en:
%  - hvilke service-knapper og animationsfunktioner den indeholder
%  - hvordan den startes (slider_control.launch.py el.l.)
%  - skærmbillede kan suppleres i AppendixA-foto-afsnittet

\textit{Indhold tilføjes.}

% ---------------------------------------------------------------------
\section{Stabil USB-portbinding -- konfigurationsændring}
\label{sec:appB-portbinding}

% TODO: Den faktiske konfigurationsfil eller kode-diff, der binder
% manager-noden mod /dev/serial/by-path/ frem for /dev/ttyUSBx.
% Dokumentationen er beskrevet i hovedteksten i
% kapitel~\ref{ch:el}, afsnit~\ref{sec:usb-binding};
% selve filændringen hører hjemme her.

\textit{Indhold tilføjes.}
